\documentclass[11pt]{article}

\usepackage[margin=0.5in]{geometry}
\usepackage{xcolor}
\usepackage{enumitem}
\usepackage{cite}
\renewcommand{\citepunct}{,}
\usepackage{hyperref}
\usepackage{amsmath,amssymb}
\usepackage{amsthm}

\newcommand{\blue}[1]{\textcolor{blue}{#1}}
\newcommand{\red}[1]{\textcolor{red}{#1}}
\newcommand{\rigorousstationarity}{at least a rigorous stationarity relationship}


\newtheorem{proposition}{Proposition}
\newtheorem{example}{Example}

\hypersetup{
  colorlinks=true,
  linkcolor=blue,
  citecolor=blue,
  urlcolor=blue
}

\title{Appeal Letter}
\author{}
\date{}
\begin{document}
\maketitle
\noindent Dear Editor-in-Chief,

We are writing to respectfully ask for your reconsideration of the decision on our manuscript \textbf{[Submission ID]}, ``DualHash: A Stochastic Primal-Dual Algorithm with Theoretical Guarantee for Deep Hashing.''

The central contribution of this work is to address the
fundamental quantization difficulty in deep hashing through
Fenchel duality, accompanied by rigorous oracle complexity
upper bounds, where prior deep hashing methods did not provide such theoretical guarantee. Experiments on multiple
datasets confirm the practical advantages of our framework.

We were informed of the rejection decision, along with the reviewers' comments. We have carefully considered all of the comments, and while we appreciate both the reviewers' and the Associate Editor's time and effort, we respectfully disagree with several of the points raised. We believe that some of the concerns underlying the rejection decision are based on a misunderstanding of the formulation and theoretical results presented in the manuscript. Below, we explain, point by point, the specific issues with which we disagree and provide supporting arguments.

% ============================================================
\section*{Reviewer 2, Comment 1: Stationarity Relationship}

\noindent\textbullet\quad\textbf{Reviewer comment.} \emph{%
The key question is whether this transformed problem provides meaningful guarantees for the original nonconvex W-regularized hashing objective. I do not find this connection convincingly established. In particular, the paper should clarify whether there is \blue{strong duality, zero duality gap, or \rigorousstationarity between the transformed primal--dual problem and the original objective.}}

\vspace{0.5em}

\noindent\textbullet\quad\textbf{\red{Our response.} }

We thank the reviewer for raising this point. The stationarity relationship requested by the reviewer is already built into our manuscript: it is encoded in the \emph{second optimality condition of (5.2) in Lemma~5.2}, namely the Fenchel--Young equality
\[
h(B^*)+h^*(\Lambda^*)=\langle \Lambda^*,B^*\rangle ,
\]
which Lemma~5.2 derives from the primal--dual optimality characterization in~\cite[Theorem~2.158]{bonnans2013perturbation} (reference~[6] of the manuscript). Below we make explicit what this condition implies for the original objective. The argument uses only Lemma~5.2, the closed form of $h^*$ given on p.~8 of the manuscript, and textbook facts of convex analysis~\cite[Theorem~23.5]{rockafellar1970convex}, \cite[Theorem~10.1]{rockafellar1998variational}; we did not spell it out in the original submission because it follows from these standard results in a few lines. We will add it to the revision as a corollary of Lemma~5.2.

\vspace{0.3em}
\textbf{(1) Under the Fenchel--Young equality of Lemma~5.2, every critical point of the Lagrangian (4.2) is a stationary point of Problem~(P); this equality is equivalent to the pointwise contact condition $|\widetilde{B}_{ij}|\ge 1$.}

We keep the notation of the manuscript. Problem~(P) is $\min_{x,B} J(x,B)=F(x,B)+h(B)$ with the W-type regularizer $h(B)=\lambda\sum_{i,j}\bigl||B_{ij}|-1\bigr|$, and the Lagrangian is $\mathcal{L}(x,B,\Lambda)=F(x,B)+\langle B,\Lambda\rangle-h^*(\Lambda)$, see (4.2). Conjugating the closed form of $h^*$ on p.~8 once more gives the biconjugate
\[
h^{**}(B)=\lambda\sum_{i,j}\bigl(|B_{ij}|-1\bigr)_+ ,
\]
which coincides with $h$ exactly on $\{|B_{ij}|\ge 1\ \forall i,j\}$.

\begin{proposition}[Stationarity transfer, corollary of Lemma~5.2]\label{prop:transfer}
Let $(\tilde{x},\widetilde{B},\widetilde{\Lambda})\in\operatorname{crit}\mathcal{L}$, i.e.,
\[
\tag{D-stat}\label{eq:pd-stationarity}
\nabla_x F(\tilde{x},\widetilde{B})=0,\qquad
\nabla_B F(\tilde{x},\widetilde{B})+\widetilde{\Lambda}=0,\qquad
\widetilde{B}\in\partial h^*(\widetilde{\Lambda}).
\]
Then the following statements are equivalent:
\begin{enumerate}[label=(\alph*)]
\item $(\tilde{x},\widetilde{B})$ is a first-order stationary point of Problem~(P):
\[
\tag{P-stat}\label{eq:original-stationarity}
\nabla_x F(\tilde{x},\widetilde{B})=0,\qquad
0\in\nabla_B F(\tilde{x},\widetilde{B})+\partial h(\widetilde{B});
\]
\item the Fenchel--Young equality in (5.2) holds at $(\widetilde{B},\widetilde{\Lambda})$: $h(\widetilde{B})+h^*(\widetilde{\Lambda})=\langle\widetilde{\Lambda},\widetilde{B}\rangle$;
\item the contact condition holds at $\widetilde{B}$: $h(\widetilde{B})=h^{**}(\widetilde{B})$, equivalently $|\widetilde{B}_{ij}|\ge 1$ for all $i,j$.
\end{enumerate}
\end{proposition}

The contact condition is required \emph{only at the point $\widetilde{B}$}, not on the whole space. A global identity $h=h^{**}$ would force $h$ to be convex, which the W-type regularizer is not; the reformulation is useful precisely because $h$ is nonconvex while $h^*$ is convex. The proof is given in the Appendix.

\vspace{0.3em}
\textbf{(2) The condition in Proposition~\ref{prop:transfer} cannot be dropped: without it, stationarity transfer fails and a positive duality gap arises.}

\begin{example}
Consider the one-dimensional problem
\[
\min_{b\in\mathbb{R}}\Bigl\{\tfrac12 b^2+\bigl||b|-1\bigr|\Bigr\},
\]
i.e., $F(b)=\tfrac12 b^2$ and $h(b)=||b|-1|$ with $\lambda=1$. The primal optimum is attained at $b^*=\pm1$ with $p^*=\tfrac12$. Since $\min_b\{\tfrac12 b^2+\Lambda b\}=-\tfrac12\Lambda^2$ and $h^*(\Lambda)=|\Lambda|$ on $[-1,1]$, the dual optimal value is
\[
d^*=\sup_{|\Lambda|\le1}\Bigl\{-\tfrac12\Lambda^2-|\Lambda|\Bigr\}=0,
\]
attained at $\Lambda^*=0$. Hence $p^*-d^*=\tfrac12>0$ and strong duality fails.

Moreover, the pair $(\bar b,\bar\Lambda)=(0,0)$ satisfies \eqref{eq:pd-stationarity}:
$F'(0)+\bar\Lambda=0$ and $\bar b=0\in\partial h^*(0)=[-1,1]$.
However, $0\notin\partial h(0)=\{-1,1\}$, so $\bar b$ is not stationary for the original objective. Consistently with Proposition~\ref{prop:transfer}, both the Fenchel--Young equality ($h(0)+h^*(0)=1\neq0$) and the contact condition ($|\bar b|=0<1$) fail at this point.
\end{example}

\noindent Hence the equivalence in Proposition~\ref{prop:transfer} is sharp: the contact condition marks exactly where stationarity transfer holds for the W-type regularizer.

\vspace{0.3em}
\textbf{(3) Strong duality cannot be expected here; the appropriate guarantee is a stationarity relationship, which Lemma~5.2 and Proposition~\ref{prop:transfer} provide.}

The counterexample in~(2) rules out both zero duality gap and unconditional stationarity transfer for the W-type regularizer. This is consistent with the nonconvex primal--dual literature, where Lagrangian-based methods are analyzed without assuming strong duality and the relationship between the formulations is established at the level of stationarity~\cite{Bolte2018,li2015global}. Our analysis follows the same line: Theorems~5.2 and~5.3 bound the SFO complexity of reaching $\varepsilon$-critical points of $\mathcal{L}$, and Proposition~\ref{prop:transfer} states precisely when such points are stationary for Problem~(P).

\section*{Reviewer 3, Comments: Clarification of Assumptions and Corrections to the Theoretical Analysis}

\noindent\textbullet\quad\textbf{Reviewer comment.} \emph{%
The detailed review of the convergence and complexity analysis identifies several substantive issues, including insufficient assumptions, incomplete boundedness and descent arguments, an underdeveloped stochastic framework, and errors in conditional expectations and estimates. Some of these issues affect the main complexity results, so the current theoretical guarantees are not fully established as written.}
 
\noindent\textbullet\quad\textbf{\red{Our response.}}

We thank Reviewer~3 for the careful examination of our theoretical analysis. Below, we clarify several major concerns and explain how we will address the remaining issues.

\vspace{0.5em}

\noindent\textbf{(1) Clarification of the assumptions
already stated in the manuscript.}

The nonemptiness of the stationary-point set was stated in the original manuscript on p.~[10]:
For this specific assumption, we respectfully clarify that it was not omitted. We will make its role more
explicit in the revision.

\vspace{0.5em}

\noindent\textbf{(2) Explicit justification of the
boundedness of the objective and descent arguments.}

Although our analysis already establishes the boundedness properties of the constructed Lyapunov function, we agree that the relevant conditions should be stated more explicitly. Under the \red{additional condition used in our analysis}, the objective function is also bounded. We will clarify these assumptions and their roles in the revised manuscript. 

The boundedness condition is crucial for ensuring reliable
convergence of the iterative sequence in stochastic con-
strained optimization. The inherent randomness of the
process makes it difficult to guarantee that all iterates re-
main within a specific level set. The necessity of
Assumption~5.1(i) is well established in
\cite{jin2022stochastic,shi2025momentum,jin2025stochastic,
bian2026stochastic}.

\vspace{0.5em}

\noindent\textbf{(3) Corrections to conditional expectations
and estimates, with their implications for the main results.}

% We will define the probability space and filtration,
% specify the conditioning sigma-algebra in each relevant
% calculation, and use consistent conditional-expectation
% notation throughout Section~5 and Appendices~C and~D.
% We will also specify which statements hold almost surely.

\red{[To be completed: present the corrected inequalities
and trace their consequences through the affected
lemmas and Theorems~5.2 and~5.3. State precisely whether
the corrections affect constants, parameter conditions,
assumptions, or complexity orders.]}

\section*{}

In brief, we have provided clarifications for the
primary theoretical concerns raised in the decision
letter. We also note that Reviewer~3 recommended
major revision. We respectfully ask
the Editor-in-Chief to reconsider the decision. 

\vspace{1.5em}

\noindent Sincerely,\\
Authors of \textbf{Submission [ID]}. 

\bibliographystyle{siamplain}
\bibliography{revision_refs}

\clearpage
\appendix
\section*{Appendix: Proof of Proposition~\ref{prop:transfer}}

\begin{proof}
\emph{Step 0: a consequence of \eqref{eq:pd-stationarity}.}
Since $h^{**}$ is proper, closed and convex, and $(h^{**})^*=h^*$, the subgradient inversion rule~\cite[Theorem~23.5]{rockafellar1970convex} gives
\[
\widetilde{B}\in\partial h^*(\widetilde{\Lambda})
\iff
\widetilde{\Lambda}\in\partial h^{**}(\widetilde{B})
\iff
h^{**}(\widetilde{B})+h^*(\widetilde{\Lambda})=\langle\widetilde{\Lambda},\widetilde{B}\rangle .
\tag{A.1}
\]
In addition, \eqref{eq:pd-stationarity} gives $\widetilde{\Lambda}=-\nabla_B F(\tilde{x},\widetilde{B})$, so \eqref{eq:original-stationarity} holds if and only if $\widetilde{\Lambda}\in\partial h(\widetilde{B})$.

\emph{(b)$\iff$(c).} By (A.1), the Fenchel--Young equality for $h$ holds at $(\widetilde{B},\widetilde{\Lambda})$ if and only if $h(\widetilde{B})=h^{**}(\widetilde{B})$. Comparing $h$ with $h^{**}(B)=\lambda\sum_{i,j}(|B_{ij}|-1)_+$ componentwise, this is equivalent to $|\widetilde{B}_{ij}|\ge1$ for all $i,j$.

\emph{(b)$\Rightarrow$(a).} The equality $h(\widetilde{B})+h^*(\widetilde{\Lambda})=\langle\widetilde{\Lambda},\widetilde{B}\rangle$ together with the definition of $h^*$ means that $\widetilde{B}$ globally minimizes $h(\cdot)-\langle\widetilde{\Lambda},\cdot\rangle$. By Fermat's rule~\cite[Theorem~10.1]{rockafellar1998variational} and Proposition~A.2 of the manuscript, $\widetilde{\Lambda}\in\widehat{\partial}h(\widetilde{B})\subseteq\partial h(\widetilde{B})$, which is \eqref{eq:original-stationarity}.

\emph{(a)$\Rightarrow$(c).} By Step~0, (a) means $\widetilde{\Lambda}\in\partial h(\widetilde{B})\cap\partial h^{**}(\widetilde{B})$. Both functions are separable, so it suffices to examine a scalar component $b=\widetilde{B}_{ij}$. If $|b|<1$, then $\partial h^{**}(b)=\{0\}$, whereas $\partial h(b)=\{-\lambda\,\mathrm{sgn}(b)\}$ for $b\neq0$ and $\partial h(0)=\{-\lambda,\lambda\}$; in all cases $0\notin\partial h(b)$, so the intersection is empty. Hence $|\widetilde{B}_{ij}|\ge1$ for all $i,j$.
\end{proof}

% Bibliography entries to add to revision_refs.bib:
% bonnans2013perturbation : J. F. Bonnans and A. Shapiro, Perturbation Analysis of Optimization Problems, Springer, 2000; reprint 2013 (ref. [6] of the manuscript)
% rockafellar1970convex   : R. T. Rockafellar, Convex Analysis, Princeton University Press, 1970
% rockafellar1998variational : R. T. Rockafellar and R. J.-B. Wets, Variational Analysis, Springer, 1998 (ref. [47] of the manuscript)
\end{document}